problem:  
Let \( N_{k,l} = \mathrm{SU}(3)/\mathrm{U}(1)_{k,l} \) be an Aloff–Wallach space with the standard decomposition \( \mathfrak{m} = \mathfrak{m}_1 \oplus \mathfrak{m}_2 \oplus \mathfrak{m}_3 \) into irreducible \(\mathrm{Ad}(\mathrm{U}(1)_{k,l})\)-modules. Every \(\mathrm{SU}(3)\)-invariant metric is determined by positive parameters \((\lambda_1,\lambda_2,\lambda_3)\).  
Under the normalized homogeneous Ricci flow,
\[
\frac{d g_t}{d t} = -2\,\mathrm{Ric}(g_t) + \frac{2}{7}\mathrm{Scal}(g_t)\,g_t,
\quad g_0 = g_{\lambda_1,\lambda_2,\lambda_3},
\]
the metric evolution reduces to an autonomous system on \( \mathbb{R}_{>0}^3 / \mathbb{R}_{>0} \).  

**Problem:**  
1. Determine whether, for some pairs \((k,l)\), the normalized invariant Ricci flow on \(N_{k,l}\) admits a **nontrivial periodic orbit** in the parameter space \( \mathbb{R}_{>0}^3 / \mathbb{R}_{>0} \)—i.e., a closed trajectory distinct from fixed points corresponding to Einstein metrics.  

2. If such orbits exist, characterize the corresponding oscillatory behavior in the sectional curvatures \(K_{ij}(t)\) of invariant planes \( \mathfrak{m}_i \wedge \mathfrak{m}_j \), and determine whether their time averages coincide with the values at any Einstein metric.  

3. If no periodic orbits exist, investigate whether the normalized flow is **strictly gradient-like**, meaning there exists a smooth function \(F(\lambda_1,\lambda_2,\lambda_3)\) (e.g. a multiple of the scalar curvature functional) whose value decreases monotonically along all nonstationary invariant flows on each \(N_{k,l}\).  

Why is it a "good" problem:  
This problem examines a fundamental but unresolved feature of Ricci flow dynamics on homogeneous spaces: whether **periodic geometric evolutions** (limit cycles) can occur on compact seven-dimensional homogeneous manifolds like the Aloff–Wallach spaces. It challenges the prevailing expectation that homogeneous Ricci flows behave as gradient systems, converging only to fixed Einstein metrics.  
By seeking either the existence or impossibility of such periodic orbits—linked to the algebraic structure of the isotropy weights—this question sits precisely at the intersection of **geometric analysis** and **dynamical systems theory**.  
A resolution would clarify whether the Ricci flow on these spaces inherently admits a potential functional acting as a Lyapunov function or whether richer dynamical phenomena occur. Thus, the problem has a simple form yet deep implications, bridging **curvature evolution**, **homogeneous geometry**, and **algebraic dynamics**, making it a paradigmatic “narrow entrance, profound depth” research problem.